\paragraph{Conjugate-coordinate class: the supertrace at an arbitrary matrix fluctuation, and the species constants that can multiply an Einstein--Hilbert-type term (conditional; bounded no-go).}
Keep the setting of the preceding paragraph: Poisson limit, leaf \(\mathbb{CP}^{1,2}\), \(T^a=M^{a4}\), background gauge function, ghost, vector with \(\mu^2=-3\), \(n_{\mathrm{ex}}\) further scalars of mass \(\mu_i^2\), the supplied sixteen real fermion components of weight \(-\tfrac12\) with no bilinear fermion mass, and one window function \(\chi\) of the commuting operator \(\Delta_\kappa\) for all species. Let \(X^a\) be arbitrary functions on the leaf, \(D_a=\{X_a,\cdot\}\), \(F_{ab}=\{X_a,X_b\}\), \(\mathcal Z_{ab}=\{F_{ab},\cdot\}\) and \(\Box_X=\eta^{ab}D_aD_b\).

\emph{Operators at arbitrary \(X\).} The Jacobi identity gives \([D_a,D_b]=\mathcal Z_{ab}\). The second variation of the quartic term plus the bilinear term about \(X\) is \(\tfrac12\int\varepsilon^aL_a\) with \(L_a=-\Box_X\varepsilon_a+D_aD_b\varepsilon^b-2\{F_{ab},\varepsilon^b\}+\mu^2\varepsilon_a\); in the background gauge \(D_a\varepsilon^a\) the ghost operator is \(-\Box_X\), the further scalars at vanishing transverse matrices have \(-\Box_X+\mu_i^2\), and the square of the Dirac operator is \(\mathcal D^2=\Box_X+\tfrac12\Gamma^{ab}\mathcal Z_{ab}\) (these three identities were checked symbolically for random polynomial \(X\) on a four-dimensional phase space). Every species is therefore \(H_i=(-\Box_X)\otimes1+Z_i+m_i\) with \(Z_i=\Sigma_i^{ab}\otimes\mathcal Z_{ab}\), the same operator \(-\Box_X\) and the same \(\mathcal Z_{ab}\) for all species, constant masses \(m_i\) (\(\mu^2\), \(\mu_i^2\), and zero for the ghost and the fermions), and constant matrices \(\Sigma_i^{ab}\) acting on the species index (for the vector, minus the generator of the vector index, which gives \(-2\{F_{ab},\varepsilon^b\}\); for the fermions, minus \(\tfrac14[\Gamma^a,\Gamma^b]\) on four copies of the real four-component spinor, because the operator compared with \(-\Box_X\) is \(-\mathcal D^2\); the traces satisfy \(\operatorname{tr}\Sigma_D\Sigma_D=\tfrac12\operatorname{tr}\Sigma_V\Sigma_V\) for all index pairs).

\emph{Species sums.} With weights \(w_i\) and dimensions \(d_i\) (\(1\cdot4\), \(1\cdot n_{\mathrm{ex}}\), \(-2\cdot1\), \(-\tfrac12\cdot16\)) one has, as exact tensor identities checked for all index tuples, \(\sum_iw_id_i=0\) and \(\sum_iw_i\operatorname{tr}(Z_i^{a_1b_1}\cdots Z_i^{a_nb_n})=0\) for \(n=1,2,3\) (not only for the totally symmetrised products); for \(n=4\) the symmetrised sum is the \(s^4\) coefficient \((f_e^2+f_m^2)^2\) of the character of the preceding paragraph and does not vanish. Because \(\Sigma_i\) acts on the species index only, the term of any Duhamel expansion of \(\sum_iw_i\operatorname{Tr}[\chi\,\Phi(H_i)]\) in powers of \(m_i\) and \(Z_i\) is a species constant times a species-independent operator trace; no commutation of \(\chi\), \(\mathcal Z_{ab}\) and \(\Box_X\) is needed. Hence, for every \(X\) and every \(\Phi\), the terms without masses and with at most three insertions of \(\mathcal Z_{ab}\) cancel identically (a check with random non-commuting matrices confirms this), and the term linear in the masses and without insertions of \(\mathcal Z_{ab}\) is \(+M_{\mathrm{str}}\operatorname{Tr}[\chi\,\Phi'(H_X)]\) (for \(\Phi=e^{-(\cdot)}\) this is \(-M_{\mathrm{str}}\operatorname{Tr}[\chi\,e^{-H_X}]\)) with \(M_{\mathrm{str}}=\sum_iw_id_im_i=4\mu^2+n_{\mathrm{ex}}\mu_i^2\). The other constants that survive at the next order are \(\sum_iw_id_im_i^2\), \(\sum_iw_im_i\operatorname{tr}\Sigma_i^{ab}\Sigma_i^{cd}\) (the vector mass alone) and the quartic constant.

\emph{Consequences.} (i) The weight-two coefficient of the preceding paragraph is, for every \(X\) and not only at \(X=T\), the mass-supertrace term: no count term \(\propto R\), no spin term quadratic in \(\mathcal Z_{ab}\) and no term cubic in it appears, whatever the dependence of \(\Box_X\) and \(\mathcal Z_{ab}\) on \(A\). (ii) At the leading symbol that term is proportional to \(M_{\mathrm{str}}\int\chi\,(\sigma_X-\mathrm i0)^{-1}\) with \(\sigma_X=\eta_{ab}\,\xi(D^a)\xi(D^b)\), which depends on \(X\) through \(\mathrm dX^a\) only. For \(A^a=\alpha^a\cos\phi\), \(\phi=q\cdot y\), with constant \(\alpha^a\), the Hamiltonian field of \(X^a\) is \(W^a-\epsilon\sin\phi\,\alpha^aY_q\), where \(W^a\) is that of \(T^a\) and \(Y_q=q_\mu\{y^\mu,\cdot\}\) (chain rule; the first-order change of the field of \(F_{ab}\) was checked through the Poisson algebra for the linear ansatz \(A^a=\alpha^a\,q\cdot y\)). This leading-symbol term depends on \(A\) only through its first derivatives, so its second variation contains no second derivative of the fluctuation and no term of Einstein--Hilbert form; the next-to-leading symbol, which involves the second jet of \(X\), is not computed. With the Gaussian window of \(\Delta_{3/4}\) the Hessian of the real (principal-value) part of this moment along \(\alpha\), at frozen coefficients and \(\phi=\pi/2\), was evaluated numerically in \(59\) random draws of orbit point and \(q\): it is indefinite in \(40\), positive definite in \(5\) and negative definite in \(14\), and along \(\alpha^a\propto q^a\) it is numerically nonzero in all draws with non-null \(q\). (iii) A term of Einstein--Hilbert form can therefore enter only through the constants listed above. The one proportional to \(M_{\mathrm{str}}\) is the mass-insertion term of the weighted \(a_2\) in the Euclidean class (there a shift of the operator of species \(i\) by \(m_i\) changes the weighted \(a_2\) by \(-m_ia_1^{(i)}\), with curvature part \(-\tfrac16M_{\mathrm{str}}\,\rho R[g_\rho]\) and connection part \(-(\sum_iw_i\kappa_im_i)\rho^2X\), \(\kappa_i\) being the weight of \(\rho^2X\) of the species; the second constant above is the vector mass times \(\kappa_V\)); its value in the six-dimensional leaf requires the next-to-leading symbol and is not computed. The pure flux constant is the order-\(F^4\) term: at the phase \(\phi=0\), where the first-jet terms vanish, and at the leading order in \(q\), the flux in the leading symbol is shifted as \(F_{ab}(\xi)\to F^0_{ab}(\xi)+\epsilon\,c\,(\xi\cdot Y_q)(q_a\alpha_b-q_b\alpha_a)\), \(c=-s\,y^4\cos\phi\), and at fixed proper time and \(\kappa=\tfrac34\) the Gaussian averages in the integrand depend on \(q\cdot y\) at fixed \(q\cdot q\), through the vector \(y^\mu\) and the window (at an orbit point with \(y^i=0\), \(\langle(\xi\cdot Y_q)^2\rangle\) changes by a factor above \(2\) over a boost of rapidity \(1.2\) of a unit \(q\)); its integrand is therefore not Lorentz invariant (the window and the leaf both single out \(y^\mu\), and this diagnostic does not separate them). (iv) No finite coefficient of this term follows in the present window: at four orbit points the integrand of the order-\(F^4\) term at \(A=0\) grows as \(s^2\) at large proper time (log-slope \(2.00\)), because the flux does not vanish on the kernel plane of \(\sigma(\Box)\); the expansion in \(\mathcal Z_{ab}\) is not the expansion in a small parameter there, and a sign of an Einstein--Hilbert-type coefficient does not follow. The non-expanded leading-symbol logarithm is \(\sum_iw_i\ln\det(\sigma+Z_i)=\ln(\sigma^2-4f_m^2)+\ln(\sigma^2+4f_e^2)+4\ln\sigma-2\ln[(\sigma+f_m)^2+f_e^2]-2\ln[(\sigma-f_m)^2+f_e^2]\), whose large-\(\sigma\) series starts with \(-6(f_e^2+f_m^2)^2/\sigma^4\) (checked); its second variation under \(A\) is not computed.

This is a statement for the Poisson limit, one loop, the commuting-symbol (leading) level, the four-matrix class with \(\mu^2=-3\), the supplied fermion content with no bilinear fermion mass and the same window for all species, applied term by term in an expansion that is not shown to converge. It does not compute the next-to-leading symbol, the dependence on \(A\) of the terms of Einstein--Hilbert form, the case of transverse matrices with flux or of species-dependent windows, Moyal terms, or a window built from dynamical matrices; it identifies no induced metric with the emergent metric of Chapter~\ref{ch:vierbein-emergence} and gives no value or sign of \(G_{\mathrm{eff}}\). This changes no ledger label; Problem~F remains open.
